% ==============================================================================
% MÓDULO CAPÍTULO II: DINÁMICA SOCIOECONÓMICA Y ASEGURAMIENTO PÚBLICO
% Archivo: cap_2/2.4_socioeconomia_sisfoh.tex
% Enfoque: 100% Calidad de Atención (Donabedian y Watson)
% ==============================================================================
\subsection{Dinámica socioeconómica, vulnerabilidad y nivel de aseguramiento público}
\label{sec:2.4_economia}

La actividad socioeconómica de los hogares usuarios del Centro de Salud Belén se inserta de manera mayoritaria en el régimen de la economía popular informal y de subsistencia diaria. La Población Económicamente Activa (PEA) labora preponderantemente en el comercio ambulatorio de verduras, frutas y alimentos preparados en el Mercado Central de Belén y su concurrida feria sabatina; en el transporte menor mediante la conducción de mototaxis; en faenas eventuales de construcción civil (albañilería y peonaje); en talleres artesanales independientes; y en el trabajo doméstico no remunerado o remunerado por jornada.

De acuerdo con los registros del Sistema de Focalización de Hogares (SISFOH), el \textbf{68.2\% de las familias del ámbito de influencia se encuentra clasificada en situación de pobreza y pobreza extrema}. En coherencia con este perfil de precariedad material, el \textbf{94.5\% de la población usuaria se encuentra afiliada al régimen subsidiado del Seguro Integral de Salud (SIS)}.

Esta dependencia casi total del aseguramiento público estatal confiere una trascendencia vital a la gratuidad y suficiencia del servicio. Los usuarios acuden al centro asistencial bajo la certeza legítima de que el Estado proveerá tanto el acto clínico humanizado como los medicamentos necesarios para su recuperación. Cuando la farmacia del establecimiento incurre en roturas de stock o desabastecimiento de medicamentos esenciales (como antibióticos de primera línea, antipiréticos, micronutrientes para la anemia o fármacos antihipertensivos), se suscita una crisis devastadora en la economía doméstica: el usuario se ve conminado a realizar un \textbf{gasto de bolsillo imprevisto} en las boticas y farmacias comerciales privadas que proliferan en los alrededores del centro de salud.

Esta vulnerabilidad socioeconómica agudiza la sensibilidad del usuario frente a la calidad del trato interpersonal: cuando el paciente percibe indiferencia, respuestas displicentes o discriminación por su condición humilde al momento de solicitar atención o retirar sus medicinas, la vivencia de calidad se fractura profundamente, generando sentimientos de desamparo frente a la precariedad del sistema público.
